\section{Discussion}
\label{sec:discussion}

\paragraph{What the model actually does.}
Reading the replies, three strategies account for nearly all of the successes. (i)~\emph{Analytic solution.} Where the loss has a short closed form, the model derives it and solves it: for H$_2$ it writes $E(\theta) = -1.0520 - 0.7782\cos\theta + 0.1820\sin\theta$ and returns $\theta^\ast = -\arctan(0.1820/0.7782)$, which is the exact minimizer; for QAOA at $p = 1$ it writes the expectation of the cut as a trigonometric polynomial in $(\gamma, \beta)$, triangles included, and maximizes it. (ii)~\emph{Physical interpolation.} For the Ising chain it reasons about the two limits (all spins aligned at $h \to 0$, all spins along $X$ at $h \to \infty$) and places the angles between them according to $h$; its weights for the field term grow with $h$ and those for the coupling shrink. (iii)~\emph{Geometric feature design.} For digits it names regions of the image (``middle-left ink'', ``bottom closure'', ``upper loop''), assigns each region a sign and a magnitude, and sums them in a rotation. All three are forms of reasoning about what a weight does; none is search.

\paragraph{Why the free-design circuits are linear classifiers.}
Consecutive rotations about one axis add their angles, so a chain $ry(w_1 x_{i_1})\,ry(w_2 x_{i_2}) \cdots ry(b)$ on one qubit prepares a state whose probability of reading 1 is $\sin^2\!\big((b + \sum_k w_k x_{i_k})/2\big)$: a linear score passed through a smooth threshold. This is the structure the model found in every unconstrained trial, and it is the natural thing to find, because it is the only circuit in the language whose weights have the meaning ``how much does this pixel count''. The structured condition, which forbids the chain, shows that the model can also program two-qubit circuits, but it does so by reproducing the same logic across qubits and combining the results with controlled rotations. We see this as the central lesson of the free-design experiments: the model is good at programming weights whose effect it can state in a sentence, and it chooses architectures accordingly.

\paragraph{Scaling.}
The fixed-ansatz results (Figure~\ref{fig:scaling}) show the limit cleanly. The amplitude-encoded circuit is a fixed unitary applied to the image vector, so each class score is the squared projection of the image onto one direction, and the weights control that direction only through a product of many rotations. At 12--18 weights the model can still say what the circuit computes (``row-band contrast'', in its words) and its weights beat the optimizer; at 30--54 weights it cannot, it says so (``hand-designed heuristic weights''), and its accuracy decays toward chance while the optimizer's improves. For ten classes the circuit needs at least 54 weights before an optimizer makes it useful, and the model is at chance throughout. Weight programming therefore scales with the \emph{interpretability} of the weight space rather than with its size: the HVA has only four to eight weights, each with a physical meaning, and the model handles them well; the 30-weight HEA for the critical Ising chain is also handled (ratio 0.999), but there the model treats most weights as zero and programs a few. Problems whose solution needs many weights with no individual meaning are beyond the method, and the model knows it. The contrast with \citet{takehi2026sparks}, where the same model wrote networks of $10^4$--$10^6$ weights, is instructive: those weights live in architectures the model can decompose (a convolution that detects a paddle, a layer that reads a board), so a million numbers can be written as a few hundred decisions. Fifty-four angles whose only effect is through a product of rotations on a $64$-dimensional amplitude vector offer no such decomposition, and are harder to program than a million structured ones.

\paragraph{Warm starts.}
Even where the model's weights are not optimal, they are a good place to start. Across all VQE and QAOA trials the optimizer started from the model's angles reached its target after a median of one loss evaluation, that is, the model's point was already within tolerance, and in the remaining cases after tens rather than hundreds of evaluations. Since loss evaluations are runs of the quantum device, this is the most direct practical consequence of the results: a few seconds of model reasoning replaces most of the optimization loop.

\paragraph{Conventions and failure modes.}
The one systematic failure is a sign. In six QAOA trials every mixing angle came out negated; the magnitudes were right and the sign-corrected circuits match the optimizer. The prompt stated the convention $U_B(\beta) = e^{-i\beta\sum X}$ explicitly, and the gate list showed $rx(2\theta)$, so this is a reasoning slip rather than missing information, but it illustrates that weight programming is brittle in a way optimization is not: an optimizer does not care which direction is positive. In practice a single loss evaluation would catch it; we chose not to allow one, to keep the test clean.

\paragraph{Limitations.}
All circuits are small (at most eight qubits) and simulated without noise. The problems are well known, and the model has certainly seen the Ising model, QAOA on regular graphs and the digits data set during training; part of what we measure is recall. We regard this as a feature of the method rather than a confound (a programmer's prior knowledge is the point), but it means the results do not predict performance on novel Hamiltonians or data sets. The optimizer baseline is deliberately modest (L-BFGS-B with finite differences, a few restarts, a 60-second budget) and is a reference, not a tuned upper bound; this matters most for the deep fixed-ansatz QML cases, where a stronger optimizer would widen the gap to the model. We report one model; a second frontier model (Claude Fable 5.1) ran a subset of the experiments and its traces are released with the others, but it was not run to completion and is not analysed here. Finally, our scoring never reaches the model, which is the strictest possible setting; allowing even one round of feedback would likely remove the sign failures and improve the deep circuits, and is the obvious next experiment.
